% Generated draft: requires mathematical and source review.
\documentclass[11pt]{article}
\usepackage[margin=25mm]{geometry}
\usepackage{fontspec}
\setmainfont{Latin Modern Roman}
\usepackage{amsmath,amssymb,mathtools,unicode-math}
\setmathfont{Latin Modern Math}
\usepackage{enumitem,needspace,xurl,hyperref}
\setlength{\parindent}{0pt}
\setlength{\parskip}{5pt}
\setlength{\emergencystretch}{3em}
\newcommand{\N}{\mathbb{N}}
\newcommand{\Z}{\mathbb{Z}}
\newcommand{\Q}{\mathbb{Q}}
\newcommand{\R}{\mathbb{R}}
\newcommand{\C}{\mathbb{C}}
\newcommand{\F}{\mathbb{F}}
\newcommand{\A}{\mathbb{A}}
\newcommand{\PP}{\mathbb{P}}
\newcommand{\E}{\mathbb{E}}
\newcommand{\eps}{\varepsilon}
\newcommand{\dd}{\,\mathrm{d}}
\begin{document}
\section*{UnsolvedMath ID: 30002203}
\textbf{Source problem code:} catalogue entry
\subsection{Definitions and mathematical statement}
Let \(\mathcal A\) denote the arrangement called the ``pentagon hyperplane arrangement'' in the supplied record, and let \(M\) denote its complement. The record does not specify the hyperplanes, ambient space, or particular combinatorial symmetry. Write \(\alpha\) for the homology automorphism attributed to that symmetry. Its coefficient system and precise action are also unspecified.

The supplied problem asks two questions:
\begin{enumerate}
\item Does \(\alpha\) lift to an automorphism of the \(E^{\infty}\)-coalgebra associated with the complement?
\item Does \(\alpha\) lift to an automorphism of the fundamental group of the complement?
\end{enumerate}

A conditional precise formulation of the first question is as follows. If the intended coalgebra is a specified chain complex \(C\) with the required structure, coefficients are specified as \(R\), and an identification \(\theta\colon H_*(C)\xrightarrow{\sim}H_*(M;R)\) is given, ask whether there exists an automorphism \(F\) of that coalgebra such that
\[
\theta\circ H_*(F)\circ\theta^{-1}=\alpha.
\]
Here \(C\), \(R\), \(\theta\), and the meaning of automorphism are placeholders for missing source specifications, not additional universally quantified data. The record uses \(E^{\infty}\); whether this means an \(E_\infty\)-coalgebra, and which model and notion of equivalence are intended, requires clarification.

For the second question, after specifying a basepoint \(m\), the proposed lift would be an automorphism of \(\pi_1(M,m)\). The record does not specify how such an automorphism is to induce the given automorphism of the entire complement homology. For the path component \(M_m\) containing \(m\), there is a canonical identification
\[
\pi_1(M,m)^{\mathrm{ab}}\cong H_1(M_m;\Z).
\]
Thus, if integral degree-one homology is the intended target, a precise lifting condition would be \(\varphi^{\mathrm{ab}}=\alpha_1\) under this identification. This conditional degree-one formulation does not resolve the stated question about complement homology in all degrees. No hypothesis identifying the homology of the complement with group homology is supplied. Consequently, the full second lifting condition remains unspecified.
\subsection{Short English statement}
The problem asks whether a particular symmetry-induced homology automorphism of the pentagon arrangement complement lifts to its \(E^{\infty}\)-coalgebra and to its fundamental group. The arrangement, symmetry, coefficients, coalgebra model, and meaning of the fundamental-group lift require clarification. The supplied record labels these questions open; this draft does not independently verify that status.
\subsection{Sources}
\begin{itemize}
\item Supplied local UnsolvedMath registry record, permanent ID \(30002203\), title: ``Lifting Pentagon-Arrangement Symmetries to Homotopy Invariants.'' Exact source problem code: \texttt{catalogue entry}. Collection: Oberwolfach Reports Open Problems. No individual proposer is identified.
\item Supplied source citation: \emph{Mini-Workshop: Cohomology Rings and Fundamental Groups of Hyperplane Arrangements, Wonderful Compactifications, and Real Toric Varieties} (2013), \emph{Oberwolfach Reports}, pp.\ 2939--2983, DOI \texttt{10.4171/owr/2012/49}. Supplied source URL and literature-source link, labelled ``original OWR report'': \url{https://doi.org/10.4171/owr/2012/49}. The record describes Oberwolfach Reports as documenting mathematical workshops and their open-problem discussions.
\item Literature assessment dated exactly \texttt{2026-08-21}: the supplied curation labels the questions open and reports that no complete later resolution was identified. It lists the original OWR report as the source checked. This is supplied dated curation, not an independently verified literature assessment; time-sensitive status and cited claims require independent verification.
\item The supplied record carries the label \texttt{corrected\_verified} and states that its mathematical content was checked against the official report, with notation, context, and question form normalized without changing the source-level claim. This is a supplied verification claim; the report has not been independently consulted for this draft. The proposed year and year tag are 2012, while the supplied citation gives 2013; their relationship is not explained in the record.
\end{itemize}
\end{document}
